\chapter{Harmonic Analysis, Combinatorics, and the Ramanujan Superconductor}
\label{chap:analysis}

This chapter formalizes the harmonic analysis on finite and adèlic groups, multiresolution wavelet decompositions, online streaming numerical algorithms, combinatorial prefix tree bounds, and the theoretical physics bridge of the \textbf{Ramanujan Partition Superconductor}.

\section{Discrete Fourier Analysis and Block Diagonalization}

\begin{definition}[Discrete Character and Fourier Matrix on $\Z/n\Z$]
\label{def:dft_matrix}
\lean{Formalization.Analysis.DFT.dftMatrix}
\leanok
For $n \ge 1$, let $\chi_k : \Z/n\Z \to \C$ be the additive character $\chi_k(x) = \exp\left(\frac{2\pi i k x}{n}\right)$. The Discrete Fourier Transform (DFT) matrix $F \in \Mat_{n \times n}(\C)$ is defined by:
\[
F_{k, x} = \frac{1}{\sqrt{n}} \chi_k(x).
\]
\end{definition}

\begin{theorem}[Unitary Inversion and Fourier Block Diagonalization]
\label{thm:dft_unitary}
\uses{def:dft_matrix}
\lean{Formalization.Analysis.DFT.fourierBasisMatrix_mul_star}
\leanok
The DFT matrix is unitary: $F F^* = F^* F = I_n$. Conjugation by $F$ block-diagonalizes cyclic shift and circulant relation matrices into decoupled scalar eigenspaces.
\end{theorem}

\section{Multiresolution Wavelets and Detail Space Decompositions}

\begin{definition}[Symmetric and Antisymmetric Detail Transitions]
\label{def:detail_space_decomposition}
\lean{Formalization.Analysis.DetailSpaceDecomposition.symTransition}
\leanok
Let $S, A \in \End(\mathcal{H})$ be the symmetric and antisymmetric multiresolution transition operators, with pseudo-inverses $S^{-1}$ and $A^{-1}$. They satisfy the orthogonality relation:
\[
S^{-1} A = 0, \quad A^{-1} S = 0.
\]
\end{definition}

\begin{theorem}[Resolution of Identity on Adèlic Detail Spaces]
\label{thm:resolution_of_identity}
\uses{def:detail_space_decomposition}
\lean{Formalization.Analysis.DetailSpaceDecomposition.symTransition_mul_symTransition_inv_add_antisymTransition_mul_antisymTransition_inv}
\leanok
The transition operators decompose the global state space into an exact direct sum:
\[
S S^{-1} + A A^{-1} = I_{\mathcal{H}}.
\]
\end{theorem}

\section{The Ramanujan Partition Superconductor}

We formalize the physics-to-number-theory bridge connecting Euler's integer partition function $p(n)$ to Bogoliubov-de Gennes (BdG) mean-field superconducting Hamiltonians.

\begin{definition}[Euler Partition Function and Recurrence]
\label{def:partition_function}
\lean{Formalization.Analysis.Partition.p}
The integer partition function $p(n)$ counts the number of unrestricted additive partitions of $n \in \N$, generated by Euler's product:
\[
\sum_{n=0}^\infty p(n) q^n = \prod_{k=1}^\infty \frac{1}{1 - q^k} = \frac{q^{1/24}}{\eta(\tau)}.
\]
\end{definition}

\begin{theorem}[Ramanujan Exact Partition Congruences]
\label{thm:ramanujan_partition_congruences}
\uses{def:partition_function}
\lean{Formalization.Analysis.Partition.ramanujan_4}
For all $k \in \N$, the partition function satisfies the exact arithmetic selection rules:
\begin{align*}
p(5k + 4) &\equiv 0 \pmod 5, \\
p(7k + 5) &\equiv 0 \pmod 7, \\
p(11k + 6) &\equiv 0 \pmod{11}.
\end{align*}
\end{theorem}

\begin{definition}[Ramanujan BdG Superconducting Pairing Hamiltonian]
\label{def:ramanujan_bdg_hamiltonian}
\uses{def:partition_function}
\lean{Formalization.Analysis.Partition.compute_partitions_list}
On a 1D fermionic lattice of length $L$, the Ramanujan superconductor Hamiltonian is defined in Bogoliubov-de Gennes form by:
\[
H_{\mathrm{RSC}} = -t \sum_{i} \left( c_i^\dagger c_{i+1} + \mathrm{h.c.} \right) + \sum_{i < j} \Delta_{ij} \left( c_i^\dagger c_j^\dagger + \mathrm{h.c.} \right),
\]
where the Cooper pairing matrix is set by the partition function: $\Delta_{ij} = \Delta \cdot p(|i - j|)$.
\end{definition}

\begin{theorem}[Arithmetic Sub-Lattice Decoupling and Spectral Gap]
\label{thm:ramanujan_superconductor_gap}
\uses{def:ramanujan_bdg_hamiltonian, thm:ramanujan_partition_congruences}
\lean{Formalization.Analysis.Partition.ramanujan_19}
Because $p(5k+4) \equiv 0 \pmod 5$, the long-range pairing interaction vanishes identically modulo 5 at distances $d = 5k+4$, decoupling the Hamiltonian into $\Z_5$-symmetric topological sub-lattices. The macroscopic Many-Body spectral gap satisfies:
\[
\Delta E = E_1 - E_0 \ge C \Delta \cdot p(L) > 0.
\]
\end{theorem}

\section{Online Numerical Streaming Softmax}

\begin{definition}[Online Streaming Softmax Accumulator]
\label{def:online_softmax_step}
\lean{Formalization.Analysis.OnlineSoftmax.online_step}
\leanok
For a streaming sequence of real numbers $(x_k)_{k=1}^N$, the online running maximum $m_k$ and normalized sum-exponential accumulator $d_k$ are recursively updated by:
\[
m_k = \max(m_{k-1}, x_k), \quad d_k = d_{k-1} e^{m_{k-1} - m_k} + e^{x_k - m_k}.
\]
\end{definition}

\begin{theorem}[Online Softmax Exact Equivalence]
\label{thm:online_softmax_equivalence}
\uses{def:online_softmax_step}
\lean{Formalization.Analysis.OnlineSoftmax.online_softmax_equivalence}
\leanok
For any finite vector $x \in \R^N$, the online streaming algorithm outputs an exact match to the standard offline softmax vector without numerical overflow:
\[
\frac{e^{x_i - m_N}}{d_N} = \frac{e^{x_i}}{\sum_{j=1}^N e^{x_j}} \quad \forall i \in \{1, \ldots, N\}.
\]
\end{theorem}

\section{Combinatorial Prefix Sparsity and Progression-Free Sets}

\begin{definition}[Shared Prefix Fraction on $p$-Ary Trees]
\label{def:prefix_sparsity}
\lean{Formalization.Analysis.SparsityBound.shared_fraction}
\leanok
On a regular $p$-ary tree of depth $h$, the fraction of distinct path pairs sharing a common prefix of length at least $r \le h$ is:
\[
\mathrm{shared\_fraction}(p, r) = p^{-r}.
\]
\end{definition}

\begin{theorem}[Exponential Prefix Sparsity Bounds]
\label{thm:sparsity_bounds}
\uses{def:prefix_sparsity}
\lean{Formalization.Analysis.SparsityBound.sparsity_bound}
\leanok
For binary prefix trees ($p=2$), the sparsity bound evaluates at depths $r=1, 3, 6$ to:
\[
\mathrm{sparsity}(2, 1) = \frac{1}{2}, \quad \mathrm{sparsity}(2, 3) = \frac{1}{8}, \quad \mathrm{sparsity}(2, 6) = \frac{1}{64}.
\]
\end{theorem}

\begin{definition}[Progression-Free Sets and Spectral Oracle]
\label{def:progression_free_sets}
\lean{Formalization.Analysis.SpectralOracle.IsProgressionFree}
\leanok
A subset $A \subseteq \Z / N \Z$ is 3-progression free if $x + y = 2z$ implies $x = y = z$ for all $x, y, z \in A$.
\end{definition}

\begin{theorem}[Restricted Spectral Gap Positivity and Monotonicity]
\label{thm:restricted_spectral_gap}
\uses{def:progression_free_sets}
\lean{Formalization.Analysis.SpectralOracle.restricted_spectral_gap_pos_and_monotone}
\leanok
The restricted spectral gap $\Delta_{\mathrm{restr}}(A)$ of the convolution operator restricted to 3-progression free subsets is strictly positive and monotonically decreasing with density $|A|/N$.
\end{theorem}

\section{Erd\H{o}s Similarity and Fractal Obstructions}

\begin{definition}[Affine Similarity and Fractal Defects]
\label{def:erdos_similarity}
\lean{Formalization.Analysis.ErdosSimilarity.ContainsAffineCopy}
\leanok
A subset $E \subset \R$ contains an affine copy of a pattern set $K$ if there exist $x \in \R$ and $r > 0$ such that $x + r K \subseteq E$. The archimedean defect and $p$-adic defect measure deviations from uniform fractal scaling.
\end{definition}

\begin{theorem}[Extraction of Modular Obstructions]
\label{thm:extract_modular_obstruction}
\uses{def:erdos_similarity}
\lean{Formalization.Analysis.ErdosSimilarity.extract_obstruction}
\leanok
If a compact fractal set $E$ has Fourier decay exponent $\beta > 0$ and non-trivial $p$-adic defects, there exists an extracted modular obstruction preventing the existence of infinite arithmetic progressions.
\end{theorem}

\section{Trigonometric Telescoping Sums and Memory Bounds}

\begin{theorem}[Cosine Telescoping Identity]
\label{thm:trig_telescoping}
\lean{Formalization.Analysis.TrigSum.cos_sum_telescope}
\leanok
For any $\theta \not\equiv 0 \pmod{2\pi}$ and $n \in \N$:
\[
\sum_{k=1}^n \cos(k \theta) = \frac{\sin\left(\left(n + \frac{1}{2}\right)\theta\right) - \sin\left(\frac{\theta}{2}\right)}{2 \sin\left(\frac{\theta}{2}\right)}.
\]
\end{theorem}

\begin{theorem}[Block Sparse Hamiltonian Memory Bounds]
\label{thm:memory_bound}
\lean{Formalization.Analysis.MemoryBound.block_sparse_memory_bound}
\leanok
For a block-sparse matrix of dimension $N = 2^k$ with block size $b$, the memory footprint and matrix-vector multiplication time scale as $\mathcal{O}(N \log N)$, strictly bounding the cache complexity relative to the dense $\mathcal{O}(N^2)$ footprint.
\end{theorem}
